\documentclass[11pt]{article}

% ============================================================
%  中远海特风电塔筒配载算法 —— 算法设计与求解报告
%  说明：本文件可直接复制到 Overleaf 编译（引擎选 XeLaTeX）。
% ============================================================

\usepackage{ctex}
\usepackage{amsmath,amssymb,amsthm}
\usepackage{geometry}
\geometry{margin=2.4cm}
\usepackage{enumitem}
\usepackage{booktabs}
\usepackage{array}
\usepackage{longtable}
\usepackage{multirow}
\usepackage{xcolor}
\usepackage{listings}
\usepackage{algorithm}
\usepackage{algpseudocode}
\usepackage{hyperref}
\hypersetup{colorlinks=true,linkcolor=blue!55!black,citecolor=blue!55!black,urlcolor=blue!55!black}

% ---- 求解器日志代码块样式 ----
\definecolor{loggray}{gray}{0.96}
\lstdefinestyle{solverlog}{
  basicstyle=\ttfamily\scriptsize,
  backgroundcolor=\color{loggray},
  breaklines=true,
  frame=single,
  framesep=4pt,
  rulecolor=\color{gray!40},
  columns=fullflexible,
  keepspaces=true,
}

\newcommand{\unit}[1]{\,\mathrm{#1}}
\newtheorem{remark}{注}

\title{\textbf{中远海特风电塔筒配载算法}\\[4pt]\large 算法设计与求解报告}
\author{配载寻优智能体平台 · 算法说明}
\date{\today}

\begin{document}
\maketitle

\begin{abstract}
\noindent
本报告面向多层多舱船舶的风电塔筒配载问题，给出从\textbf{问题分析}、\textbf{策略选择}、
\textbf{数学模型}、\textbf{算法设计}、\textbf{参数整定}到\textbf{求解与结果分析}的完整技术路线。
问题被建模为以\emph{完整套数最大化}为目标的混合整数线性规划（MILP）：通过\emph{候选位离散化}
把连续的二维布置空间转化为有限的 0--1 选择，再以 Pyomo 建模、HiGHS 求解。
在算例（单一货型 T110.5-70A、6 个塔筒分段；4 个舱室、40 块 bypass 盖板、20 块舱盖板）上，
最优配载结果为 \textbf{27 套}（最优性间隙 $0\%$），布置货物 106 件，启用 bypass 盖板 17 块，
几何冲突 0 处。
\end{abstract}

\tableofcontents
\bigskip

% ============================================================
\section{问题分析}
% ============================================================

\subsection{业务背景}
中远海特承运的风电塔筒（Tower）按\textbf{套}发运，每一套由 6 个不同的分段构件组成
（\texttt{componentNo} $=1,\dots,6$），\textbf{必须成套}装船——只装上其中一部分构件不计入有效套数。
船舶为多层（Layer~1$\sim$4）、多舱（Hatch）结构：底层（Layer~1）货物坐落于舱内；
上部各层（Layer~2/3）货物需"担"在专用的 \textbf{bypass 盖板} 上；顶层（Layer~4）货物坐落于舱盖板上。
配载方案需同时满足几何排布、跨层限高、盖板承重等一系列工程约束。

\subsection{输入与输出}
\paragraph{输入。}
\begin{itemize}[leftmargin=1.8em,itemsep=2pt]
  \item \texttt{cargoData}：每个分段的长、宽、高（单位 mm）、重量（吨）与构件编号。
  \item \texttt{vesselStructure}：
    \begin{itemize}[itemsep=1pt]
      \item \texttt{bypassBoardPosition}：各 bypass 盖板的四角坐标、所属层与舱号；
      \item \texttt{hatchCoverPosition}：各舱盖板的四角坐标；
      \item \texttt{stowageFeasibleRange}：每层每舱的可配载多边形范围与限高（Layer~3 限高沿 $x$ 分段给出，Layer~4 不限高）。
    \end{itemize}
\end{itemize}

\paragraph{输出。}
总套数 \texttt{totalSetCount}、每件被选中货物的位置（四角坐标、所在层、层叠数 \texttt{tier}、
朝向 \texttt{direction}）、以及被启用的 bypass 盖板清单。

\subsection{算例规模}
本报告算例的构件参数见表~\ref{tab:cargo}；船舶结构为 4 个舱室、40 块 bypass 盖板
（Layer~2、Layer~3 各 20 块）、20 块舱盖板、共 11 个可配载区域（Layer~1 的 4 个舱 + Layer~2/3 各 3 个舱 + Layer~4 整片）。

\begin{table}[ht]
\centering
\caption{货型 T110.5-70A 的 6 个塔筒分段参数}
\label{tab:cargo}
\begin{tabular}{ccrrrr}
\toprule
\texttt{componentNo} & 名称 & 长 (mm) & 宽 (mm) & 高 (mm) & 重量 (t) \\
\midrule
1 & Tower Section 1 & 8\,910  & 4\,550 & 4\,640 & 72.6 \\
2 & Tower Section 2 & 12\,770 & 4\,550 & 4\,640 & 67.8 \\
3 & Tower Section 3 & 16\,130 & 4\,550 & 4\,640 & 68.0 \\
4 & Tower Section 4 & 21\,170 & 4\,550 & 4\,640 & 68.9 \\
5 & Tower Section 5 & 26\,210 & 4\,550 & 4\,630 & 64.1 \\
6 & Tower Section 6 & 25\,450 & 4\,540 & 4\,630 & 55.6 \\
\bottomrule
\end{tabular}
\end{table}

\subsection{约束清单（业务语义）}
\label{sec:constraints-biz}
配载需满足以下 10 类约束（详细数学形式见第~\ref{sec:model} 节）：
\begin{enumerate}[label=C\arabic*.,leftmargin=2.6em,itemsep=2pt]
  \item \textbf{方向--层叠耦合}：顺船（\texttt{direction}=1）可双层堆叠（\texttt{tier}=2）；横船（\texttt{direction}=0）只能单层。
  \item \textbf{成套约束}：必须凑齐 \texttt{componentNo} $1\sim6$，总套数为整数，不允许残套。
  \item \textbf{横向间距}：$y$ 向货物间距须大于 $100\unit{mm}$。
  \item \textbf{纵向间距}：$x$ 向货物间距须大于 $500\unit{mm}$。
  \item \textbf{四角落板}：Layer~2/3/4 货物四角须落在某块盖板（bypass 或舱盖板）范围内。
  \item \textbf{禁止头对头}：同一块盖板上两端不得头对头，端部至盖板边界的条带内不得侵入其他货物。
  \item \textbf{可配载范围与跨层限高}：足迹须在范围内；\texttt{tier}=2 时高度翻倍；下层货物若突破累计限高，则上层对应 bypass 盖板不得启用、且空间不得重叠。
  \item \textbf{bypass 数量上限}：启用的 bypass 盖板总数 $\le 18$，且与承载货物双向绑定。
  \item \textbf{不跨舱}：Layer~1/2/3 货物四角须位于同一舱室内。
  \item \textbf{均布承重}：每块盖板上货物半重之和除以盖板面积 $\le 3\unit{t/m^2}$。
\end{enumerate}

% ============================================================
\section{策略选择}
% ============================================================

\subsection{问题的组合本质}
该问题在数学上属于带复杂工程约束的\textbf{二维布局 / 装箱}问题，其核心难点有三：
（i）决策包含\emph{连续}的坐标位置与\emph{离散}的选取、朝向、层叠、盖板启用，是混合型决策；
（ii）目标"成套数"由\emph{瓶颈构件}决定（最稀缺构件的可装数量封顶了总套数），具有强耦合性；
（iii）约束跨层耦合（下层高度影响上层可用性），不可简单分层独立求解。

\subsection{候选位离散化 + 精确求解}
\label{sec:strategy}
直接对连续坐标做混合整数非线性优化既难建模又难保证最优。我们采用工业界成熟、可证最优的两段式策略：

\begin{enumerate}[leftmargin=2em,itemsep=3pt]
  \item \textbf{候选位离散化（Discretization）}：把"在何处放置某构件"的连续自由度，
  预先离散为有限个\emph{候选摆放位}（candidate placement）。每个候选位是一个完全确定的方案原子——
  固定了构件、层、舱、朝向、层叠数与精确坐标。\emph{平面落位、四角落板、跨舱、Layer~3 自身限高等纯几何约束在生成阶段即被过滤}，
  保证每个候选位\textbf{单体可行}。
  \item \textbf{0--1 精确优化（MILP）}：在候选位集合上引入二元选择变量，
  把"成套""间距""承重""跨层耦合"等\emph{关系型}约束写成线性不等式，
  用分支定界（Branch \& Bound）求\textbf{全局最优}并给出最优性证明（对偶间隙）。
\end{enumerate}

\begin{remark}
离散化的\emph{密度}由"分相（phase）"参数控制：密度越高，候选位越多，可逼近的最优解越好，但模型规模与求解耗时上升。
第~\ref{sec:param} 节说明如何整定该密度，使整数最优解稳定在几何上限。
\end{remark}

\subsection{为何选择 MILP 而非启发式}
启发式（遗传算法、模拟退火、贪心）能快速给出可行解，但\textbf{无法证明最优}，
且对"成套"这类全局耦合约束容易陷入局部最优。考虑到本问题离散化后规模可控、
且客户高度关注"是否已装到极限"，我们选择 MILP 精确求解：HiGHS 的 presolve 与割平面技术
可将数百万条配对约束压缩到可解规模，并在数分钟内给出\textbf{零间隙}最优证明（见第~\ref{sec:solve} 节）。

% ============================================================
\section{数学模型}
\label{sec:model}
% ============================================================

\subsection{集合与索引}
\begin{align*}
\mathcal{K} &= \{1,2,3,4,5,6\} && \text{构件编号集} \\
\mathcal{T} &= \{\text{T110.5-70A}\} && \text{货型集} \\
\mathcal{L} &= \{1,2,3,4\} && \text{船层集} \\
\mathcal{H} &= \{\text{Hatch1},\dots,\text{Hatch4}\} && \text{舱室集} \\
\mathcal{P} &  && \text{候选摆放位集；} \mathcal{P}_{t,k},\ \mathcal{P}^{\ell} \text{ 分别为货型 } t\text{、构件 } k\text{、层 } \ell \text{ 的子集} \\
\mathcal{B} &  && \text{bypass 盖板集；} \mathcal{B}^{\ell} \text{ 为层 } \ell \text{ 的子集（}\ell\in\{2,3\}\text{）} \\
\mathcal{C} &  && \text{舱盖板集（Layer~4）} \\
\mathcal{S}(p) & \subseteq \mathcal{B}\cup\mathcal{C} && \text{支撑候选位 } p \text{ 的盖板集合}
\end{align*}

\subsection{参数}
对每个候选位 $p\in\mathcal{P}$，下列量在生成阶段即被确定：
\begin{itemize}[leftmargin=1.8em,itemsep=2pt]
  \item $k(p)\in\mathcal{K}$、$\ell(p)\in\mathcal{L}$、$d(p)\in\{0,1\}$、$t(p)\in\{1,2\}$：构件编号、层、朝向、层叠数；
  \item 轴对齐足迹矩形 $[\,x_1^p,x_2^p\,]\times[\,y_1^p,y_2^p\,]$（朝向决定长边沿 $x$ 还是 $y$）；
  \item 自重 $\omega_p$（吨）与\textbf{半重} $\tilde\omega_p = \omega_p\, t(p)/2$（每端盖板分担的等效重量）；
  \item 堆垛总高 $h_p = h^{\text{cargo}}_{k(p)}\cdot t(p)$。
\end{itemize}
其余常量：横向最小间距 $g_y=100\unit{mm}$，纵向最小间距 $g_x=500\unit{mm}$，
最大均布强度 $\sigma = 3\unit{t/m^2}$，bypass 启用上限 $B_{\max}=18$，
盖板 $b$ 的面积 $A_b$（$\mathrm{m}^2$），层 $\ell$ 舱 $r$ 的限高 $H^{\lim}_{\ell,r}$。

\subsection{决策变量}
\begin{align}
y_p &\in \{0,1\}, & & p\in\mathcal{P} && \text{（是否选用候选位 } p\text{）}\\
z_b &\in \{0,1\}, & & b\in\mathcal{B} && \text{（是否启用 bypass 盖板 } b\text{）}\\
S   &\in \mathbb{Z}_{\ge 0}. & & && \text{（完整套数）}
\end{align}

\subsection{目标函数}
\begin{equation}
\boxed{\ \max\ S\ }
\end{equation}

\subsection{约束（关系型，进入 MILP）}

\paragraph{C2 成套约束。}
每种货型的每个构件，其计入层叠的装载量都不得少于总套数：
\begin{equation}
\sum_{p\in\mathcal{P}_{t,k}} t(p)\, y_p \ \ge\ S,
\qquad \forall\, t\in\mathcal{T},\ k\in\mathcal{K}.
\label{eq:c2}
\end{equation}
由于目标是最大化 $S$，其最优值等于各构件可装单元数的\emph{最小值}（瓶颈构件）。
若某构件完全无可行候选位，则其右端强制 $S=0$。

\paragraph{C3 / C4 / C6 配对互斥（同层同舱）。}
将几何上不能共存的两个候选位 $(p_1,p_2)$ 归入冲突集，禁止同时选用：
\begin{equation}
y_{p_1} + y_{p_2} \ \le\ 1,
\qquad \forall\, (p_1,p_2)\in \mathcal{C}_3\cup\mathcal{C}_4\cup\mathcal{C}_6 .
\label{eq:c346}
\end{equation}
冲突集的几何判据（仅在同层、同舱内枚举；Layer~4 按层枚举）：
\begin{itemize}[leftmargin=1.8em,itemsep=2pt]
  \item $\mathcal{C}_3$（横向）：$x$ 向投影严格重叠，且 $y$ 向净间距 $<g_y$；
  \item $\mathcal{C}_4$（纵向）：$y$ 向投影严格重叠，且 $x$ 向净间距 $<g_x$；
  \item $\mathcal{C}_6$（头对头）：对担在盖板上的 $p_1$，由其两端到盖板边界、沿 $p_1$ 的 $y$ 跨度形成的"禁区"
  若与 $p_2$ 的足迹重叠，则二者互斥。
\end{itemize}
其中"$x$ 向净间距"指当 $x$ 区间不重叠时两区间的最近距离 $\max(x_1^{p_1},x_1^{p_2})-\min(x_2^{p_1},x_2^{p_2})$，$y$ 向同理。

\paragraph{C7 跨层高度耦合。}
设低层候选位 $p$（$\ell(p)\in\{1,2\}$）位于舱 $r$，其堆垛高 $h_p$ 突破到上层 $L$ 的判据为累计限高比较
（如 $\ell(p)=1$ 时：$h_p>H^{\lim}_{1,r}$ 影响 $L=2$，$h_p>H^{\lim}_{1,r}+H^{\lim}_{2,r}$ 影响 $L=3$，依此类推）。
对每个被影响的上层 $L$：
\begin{align}
y_p + z_b &\ \le\ 1, & &\forall\, (p,b)\in\mathcal{C}_7^{\text{couple}}
  && \text{（占据上层空间 $\Rightarrow$ 对应 bypass 禁用）} \label{eq:c7a}\\
y_p + y_q &\ \le\ 1, & &\forall\, (p,q)\in\mathcal{C}_7^{\text{body}}
  && \text{（与上层货物足迹重叠 $\Rightarrow$ 互斥）} \label{eq:c7b}
\end{align}
其中 $(p,b)$ 取 $b\in\mathcal{B}^{L}$ 且其包围盒与 $p$ 足迹重叠；$(p,q)$ 取 $q\in\mathcal{P}^{L}$ 且其足迹与 $p$ 足迹重叠。

\paragraph{C8 bypass 启用约束。}
\begin{align}
\sum_{b\in\mathcal{B}} z_b &\ \le\ B_{\max}, \label{eq:c8a}\\
y_p &\ \le\ z_b, & &\forall\, p:\ \ell(p)\in\{2,3\},\ b\in\mathcal{S}(p)\cap\mathcal{B}, \label{eq:c8b}\\
z_b &\ \le\ \sum_{p:\, b\in\mathcal{S}(p)} y_p, & &\forall\, b\in\mathcal{B}. \label{eq:c8c}
\end{align}
式~\eqref{eq:c8b} 表示"承载货物 $\Rightarrow$ 必须启用其支撑盖板"；
式~\eqref{eq:c8c} 杜绝"无人使用却被启用"的虚假启用（若 $b$ 无任何使用者则强制 $z_b=0$）。

\paragraph{C10 均布承重约束。}
对 bypass 盖板（受 $z_b$ 门控）与舱盖板（恒可用）分别有：
\begin{align}
\sum_{p:\, b\in\mathcal{S}(p)} \tilde\omega_p\, y_p &\ \le\ \sigma\, A_b\, z_b, & &\forall\, b\in\mathcal{B}, \label{eq:c10a}\\
\sum_{p:\, c\in\mathcal{S}(p)} \tilde\omega_p\, y_p &\ \le\ \sigma\, A_c, & &\forall\, c\in\mathcal{C}. \label{eq:c10b}
\end{align}

\paragraph{变量上界（加速用有效不等式）。}
由 C2 可推出 $S$ 的一个紧上界，作为变量盒约束加入以收紧根松弛：
\begin{equation}
0\ \le\ S\ \le\ \min_{k\in\mathcal{K}}\ \sum_{p\in\mathcal{P}_k} t(p).
\label{eq:bound}
\end{equation}

\subsection{在生成阶段内化的约束（C1、C5、C7-平面、C9）}
C1（方向--层叠）、C5（四角落板）、C7 的\emph{平面落位}与\emph{Layer~3 自身限高}、C9（不跨舱）
均为\textbf{单体几何可行性}约束，已在候选位生成时强制满足（不可行的候选位根本不生成）。
因此它们\emph{不进入} MILP 的约束矩阵，从而显著降低模型规模。
求解后另有独立校验器复核全部 10 类约束（见第~\ref{sec:solve} 节，冲突数 $=0$）。

% ============================================================
\section{算法设计}
% ============================================================

\subsection{总体流水线}
求解器实现为五阶段流水线（与代码模块一一对应）：
\[
\text{载入} \rightarrow \text{候选生成} \rightarrow \text{建模} \rightarrow \text{求解} \rightarrow \text{输出/校验}.
\]

\subsection{候选位生成（离散化核心）}
对每个构件、每个朝向、每个可行区域，按"步距网格 + 分相偏移"枚举锚点，并逐一通过几何过滤：

\begin{algorithm}[ht]
\caption{分层候选位生成（要点）}
\label{alg:gen}
\begin{algorithmic}[1]
\Require 构件集、可配载范围、bypass/舱盖板、间距 $g_x,g_y$、各层分相参数
\State $\mathcal{P}\gets\varnothing$
\ForAll{构件 $c$，朝向 $d\in\{1,0\}$}
  \State 由 $d$ 决定足迹边长 $(s_x,s_y)$；\textbf{若} $d=0$ \textbf{则} 层叠集 $\gets\{1\}$ \textbf{否则} $\{1,2\}$ \Comment{内化 C1}
  \State 步距 $\Delta_x\gets s_x+g_x,\ \Delta_y\gets s_y+g_y$
  \ForAll{分相偏移 $(o_x,o_y)$}
    \ForAll{锚点 $(x,y)$ 沿网格 $x\!=\!x_{\min}\!+\!o_x\!:\!\Delta_x,\ y\!=\!y_{\min}\!+\!o_y\!:\!\Delta_y$}
      \If{足迹 $\not\subseteq$ 可配载多边形} \textbf{continue} \EndIf \Comment{内化 C7-平面、C9}
      \If{$\ell\in\{2,3,4\}$ 且 四角未全部落在盖板上} \textbf{continue} \EndIf \Comment{内化 C5}
      \If{$\ell=3$ 且 $h_c\!\cdot\! t$ 超过分段限高} \textbf{continue} \EndIf \Comment{内化 C7-L3}
      \State 记录支撑盖板 $\mathcal{S}(p)$，将候选位 $p$ 加入 $\mathcal{P}$
    \EndFor
  \EndFor
\EndFor
\State \Return $\mathcal{P}$
\end{algorithmic}
\end{algorithm}

\begin{itemize}[leftmargin=1.8em,itemsep=2pt]
  \item \textbf{Layer~1}：在舱内多边形上自由网格布点（无需盖板支撑）。
  \item \textbf{Layer~2/3}：货物两端须精确落在一块或相邻两块 bypass 盖板上（顺船沿 $x$ 跨板、横船沿 $y$），并通过分段限高检查。
  \item \textbf{Layer~4}：四角须落在舱盖板上。
\end{itemize}

\subsection{冲突对预计算}
关系型约束 \eqref{eq:c346}--\eqref{eq:c7b} 的系数全部在建模前以几何方式枚举：
\begin{itemize}[leftmargin=1.8em,itemsep=2pt]
  \item C3/C4：在每个（层，舱）分组内按 $x$（或 $y$）排序，用\textbf{扫描线}剪枝——
  一旦后续候选位的起点越过当前候选位的终点即可提前 \texttt{break}，避免 $O(n^2)$ 全枚举；
  \item C6：为每个担板候选位预计算两侧"禁区"矩形，再与同层其他候选位做矩形相交测试；
  \item C7：对每个低层候选位，按累计限高判定其影响的上层，再与该层 bypass 盖板、上层候选位做包围盒/足迹相交测试。
\end{itemize}

\subsection{建模与求解}
用 Pyomo 构建 \texttt{ConcreteModel}：变量 $y,z,S$，目标 $\max S$，依次施加 C2--C10。
求解器优先选 HiGHS（经 Pyomo APPSI 接口直连，回退顺序 GLPK $\to$ CBC），
开启求解器内置 presolve、割平面与启发式，最终由分支定界给出最优解与对偶界。

% ============================================================
\section{参数配置（智能体自动整定）}
\label{sec:param}
% ============================================================

\subsection{可调参数}
影响解质量与耗时的关键参数为\emph{离散化密度}（每层 $x$/$y$ 方向的分相数）与\emph{最小间距}：
\[
\big(\text{L1}_x,\text{L1}_y;\ \text{L23}_x,\text{L23}_y;\ \text{L4}_x,\text{L4}_y\big),
\qquad (g_x,g_y).
\]
分相数越大，锚点网格越密，候选位集合 $\mathcal{P}$ 越大，可逼近的整数最优越好，但模型规模与耗时同步增长。
间距 $g_x=500,\ g_y=100$ 由工艺规范固定，不参与寻优。

\subsection{调参智能体的整定过程}
"决策引擎"智能体以\emph{逐步加密}策略反复求解、观测\textbf{整数最优值 $S^\star$} 与\textbf{对偶上界}的收敛情况，
直至整数最优稳定、且对偶界确认其为全局最优。整定轨迹见表~\ref{tab:tuning}。

\begin{table}[ht]
\centering
\caption{离散化密度整定轨迹（智能体自动实验；最后一行为采用配置）}
\label{tab:tuning}
\small
\begin{tabular}{clrrcc}
\toprule
试验 & 分相配置 $(\text{L1};\text{L23};\text{L4})$ & 候选位数 & 对偶上界 & 整数最优 $S^\star$ & 备注 \\
\midrule
T1 & $(1,1;1,1;1,1)$ & 5\,040  & 27.8 & 24 & 网格过疏，错失装位 \\
T2 & $(2,1;2,1;2,2)$ & 9\,210  & 27.8 & 26 & 接近但未达上界 \\
T3 & $(2,2;2,2;2,2)$ & 13\,860 & 27.7 & 27 & 首次触达几何上限 \\
\textbf{T4} & $\mathbf{(3,2;2,2;2,3)}$ & \textbf{16\,629} & \textbf{27.0} & \textbf{27} & \textbf{整数最优稳定，间隙 0\%} \\
T5 & $(3,3;3,3;3,3)$ & 28\,150 & 27.0 & 27 & 解不再改善，仅耗时增加 \\
\bottomrule
\end{tabular}
\end{table}

\begin{remark}
随着密度提高，整数最优从 $24\!\to\!26\!\to\!27$ 单调改善并在 $T3$ 之后\textbf{饱和于 27}；
继续加密（$T5$）不再提升解质量、只增加规模。故\textbf{采用配置 $T4$}：
$(\text{L1}_x,\text{L1}_y;\text{L23}_x,\text{L23}_y;\text{L4}_x,\text{L4}_y)=(3,2;2,2;2,3)$，$(g_x,g_y)=(500,100)$。
该配置生成 \textbf{16\,629} 个候选位（Layer~1: 2\,142；Layer~2: 7\,563；Layer~3: 3\,591；Layer~4: 3\,333），
在可接受耗时内给出\textbf{零间隙}的 27 套最优解。
\end{remark}

% ============================================================
\section{问题求解与结果分析}
\label{sec:solve}
% ============================================================

\subsection{模型规模}
采用配置 $T4$ 时，原始 MILP 规模为：
\begin{center}
\begin{tabular}{lr}
\toprule
指标 & 数值 \\
\midrule
变量（列） & 16\,670 \\
\quad 其中二元 $y,z$ & 16\,669\ \ (=16\,629\ y + 40\ z) \\
\quad 整数 $S$ & 1 \\
约束（行） & 15\,344\,278 \\
非零元 & 30\,753\,404 \\
\bottomrule
\end{tabular}
\end{center}
约束行数高达 $1.5\times10^{7}$，绝大多数来自 C3/C4/C6/C7 的\textbf{两两配对互斥}不等式
\eqref{eq:c346}--\eqref{eq:c7b}——这是离散化方法的固有代价，但其结构高度可压缩（见下）。

\subsection{预处理（Presolve）}
HiGHS 的 presolve 通过\textbf{团（clique）合并、探测与替代}，将海量两两互斥约束压缩为少量团约束，
规模锐减约三个数量级：
\begin{center}
\begin{tabular}{lrrr}
\toprule
& 行 & 列 & 非零元 \\
\midrule
Presolve 前 & 15\,344\,278 & 16\,670 & 30\,753\,404 \\
Presolve 后 & 26\,169 & 14\,763 & 1\,211\,989 \\
\midrule
缩减量 & $-15\,318\,109$ & $-1\,907$ & $-29\,541\,415$ \\
\bottomrule
\end{tabular}
\end{center}
日志亦报告"目标函数以比例 1 取整"（\texttt{Objective function is integral with scale 1}），
意味着对偶界可向下取整以加速收敛。

\subsection{分支定界过程}
根节点 LP 松弛上界约 $27.77$，随割平面加入逐步收紧至约 $27.59$；
中心舍入（C）等启发式先得到劣质可行解（$S=1$），随后子-MIP（\texttt{L}）启发式在约 $433.6\unit{s}$
直接跳到 $S=27$；由于此时对偶界已 $<28$ 且目标整数，分支定界在第 1 个节点即关闭间隙，证明 $27$ 为全局最优。
关键日志节选：

\begin{lstlisting}[style=solverlog]
Running HiGHS 1.14.0
MIP has 15344278 rows; 16670 cols; 30753404 nonzeros; 16670 integer variables (16669 binary)
Presolving model
  ... Presolve reductions: rows 26169(-15318109); columns 14763(-1907); nonzeros 1211989(-29541415)
Objective function is integral with scale 1
Solving MIP model with: 26169 rows, 14763 cols (14762 binary, 1 integer), 1211989 nonzeros

Src  Proc. InQueue | Leaves Expl. | BestBound   BestSol      Gap | ... Time
 z       0       0        0  0.00%  inf         -0         Large       224.5s
         0       0        0  0.00%  27.76636542 -0         Large       264.2s
 C       0       0        0  0.00%  27.6788391   1      2667.88%       277.4s
 L       0       0        0  0.00%  27.58860458 27         2.18%       433.6s
         1       0        1 100.0%  27          27         0.00%       433.7s

Solving report
  Status            Optimal
  Primal bound      27
  Dual bound        27
  Gap               0% (tolerance: 0.01%)
  Timing            433.66 (HiGHS solve)
  Nodes             1
  LP iterations     101572
status=optimal  totalSetCount=27  selected=106  enabled bypass=17
\end{lstlisting}

\begin{remark}
求解器内核耗时 $433.66\unit{s}$；端到端总耗时 $976.6\unit{s}$，其中\textbf{建模与冲突对枚举}（生成 $1.5\times10^7$ 行）
占据主要时间。若需缩短总耗时，可在保证解质量的前提下适度降低离散化密度（参见表~\ref{tab:tuning} 中 $T3$）。
\end{remark}

\subsection{最优配载结果}
\begin{table}[ht]
\centering
\caption{最优配载结果汇总}
\label{tab:result}
\begin{tabular}{lr}
\toprule
指标 & 数值 \\
\midrule
完整套数 \texttt{totalSetCount} & \textbf{27} \\
最优性间隙 & $0\%$（已证全局最优）\\
布置货物总件数 & 106 \\
\quad Layer~1 / 2 / 3 / 4 & 36 / 13 / 21 / 36 \\
启用 bypass 盖板 & 17 / 18 \\
几何冲突（独立校验） & 0 \\
HiGHS 求解耗时 & $433.66\unit{s}$ \\
端到端总耗时 & $976.6\unit{s}$ \\
\bottomrule
\end{tabular}
\end{table}

\paragraph{结果解读。}
\begin{itemize}[leftmargin=1.8em,itemsep=2pt]
  \item \textbf{瓶颈与层叠}：总套数 27 由最稀缺构件封顶；顺船候选位通过 \texttt{tier}=2 双层堆叠，
  使每种构件计入层叠后的单元数均 $\ge 27$，从而满足成套约束 \eqref{eq:c2}。
  \item \textbf{盖板利用}：启用 17/18 块 bypass 盖板（$\approx 94\%$），逼近 $B_{\max}=18$ 的硬上限，
  说明上层空间已被充分利用；所有盖板的均布强度均 $\le 3\unit{t/m^2}$。
  \item \textbf{可行性}：独立几何校验器复核 C1--C10，冲突数为 0，方案可直接交付（输出 2D 分层图、3D 立体图、结果 JSON 与 CAD 图纸）。
\end{itemize}

\subsection{结论}
本文以候选位离散化将连续布局问题转化为 MILP，借助 HiGHS 的 presolve 与分支定界，
在算例上给出\textbf{经证明的全局最优}配载方案——\textbf{27 套}，间隙 $0\%$，零几何冲突。
该方法论可推广至其他货型与船型：只需更新输入数据并由调参智能体重新整定离散化密度，
即可获得同等质量的最优配载方案。

\end{document}
